\chapter{Introduction}
\label{sec:introduction}

This chapter introduces the Styrian Health Data case study, states the analytical objectives, and summarises the dataset and principal findings that are developed in subsequent chapters.

\section{Aims and Goals}
\label{sec:aims}


The project pursues three interconnected statistical objectives:
\begin{enumerate}
    \item \textbf{Characterise data quality} and prepare an analysis-ready sample through systematic cleaning, missing-value handling, and plausibility checks.
    \item \textbf{Identify determinants of blood pressure} using descriptive statistics, group comparisons, ordinary least squares (OLS) regression, and Random Forest models; evaluate whether self-reported wellbeing and clinical risk flags explain variation in systolic and diastolic blood pressure after adjustment for demographics.
    \item \textbf{Extend the analysis spatially and contextually} by linking individual records to regional gross domestic product (GDP), urban--rural typology, and neighbourhood structure; test for neighbourhood effects on cardiovascular indicators.
\end{enumerate}

Beyond these substantive goals, the project aims to compare regression and classification formulations, quantify predictive limits of available covariates, and document limitations inherent in voluntary kiosk data.

\section{Data Description}
\label{sec:data_description}

This section describes the data collection process, the variable structure, and the data quality issues.

\subsection{Collection}

The data were \textbf{not} generated by a designed experiment. They arise from an \textbf{observational, voluntary health-screening programme} at public exhibition terminals \cite{styrian_health_data}. Participants self-reported demographic and health information via interactive terminals; blood pressure was measured automatically by the kiosk equipment. Measurements were recorded at \textbf{three terminals} between \textbf{29 April 2006} and \textbf{29 October 2006}. The total number of observations in the dataset is \textbf{16,386} records with 19 attributes.

\subsection{Variable Structure}

Table~\ref{tab:variables} summarises the variable structure. Measured systolic and diastolic blood pressure (\texttt{messwert\_bp\_sys/dia}) serve as the primary outcome variables for most analyses, with self-estimated blood pressure values (\texttt{schaetzwert\_bp\_sys/dia}) available as auxiliary measures.

Independent (predictor) variables include demographics, self-reported health and risk factors, geography, and metadata such as timestamp and terminal number. External covariates merged in later stages include hourly temperature and humidity from the Open-Meteo Historical Weather API \cite{open_meteo}, regional GDP per capita from Statistics Austria \cite{statistik_atlas}, and urban--rural typology.

\begin{table}[h]
    \centering
    \caption{Overview of key variables in the Styrian Health Data (2006).}
    \label{tab:variables}
    \begin{tabular}{p{0.22\textwidth} p{0.68\textwidth}}
        \toprule
        \textbf{Category} & \textbf{Variables and Description} \\
        \midrule

        Identifiers &
        Unique record ID (\texttt{id}), timestamp of entry (\texttt{zeit}), and data collection terminal (\texttt{terminal}) \\

        \addlinespace

        Geography &
        Place of residence and administrative hierarchy: postal code, municipality (\texttt{gemeinde}), district (\texttt{bezirk}), and state (\texttt{bundesland}) \\

        \addlinespace

        Health Status &
        Self-reported health condition (\texttt{befinden}), measured on an ordinal scale from 1 (very good) to 5 (very poor) \\

        \addlinespace

        Demographics &
        Year of birth (\texttt{geburtsjahr}), gender (\texttt{geschlecht}), and derived age variable \\

        \addlinespace

        Risk Factors &
        Smoking status (\texttt{raucher}), known diabetes status (\texttt{blutzucker\_bekannt}), known cholesterol status (\texttt{cholesterin\_bekannt}), and ongoing treatment status (\texttt{in\_behandlung}) \\

        \addlinespace

        Blood Pressure &
        Estimated systolic and diastolic blood pressure (\texttt{schaetzwert\_bp\_sys/dia}), and measured systolic and diastolic blood pressure (\texttt{messwert\_bp\_sys/dia}) \\

        \bottomrule
    \end{tabular}
\end{table}

\subsection{Data Quality Issues}

Several quality issues were identified during exploratory analysis and cleaning:
\begin{itemize}
    \item \textbf{Missing values:} Column-wise null counts include 331 missing state codes, 45--56 missing estimated blood pressure values, and 23 records with joint missingness in \texttt{befinden}, birth year, gender, and age.
    \item \textbf{Geographic representativeness:} Approximately 89\% of records with state information originate from Steiermark; the sample is not nationally representative.
\end{itemize}

\section{Main Results at a Glance}
\label{sec:results_preview}

\begin{itemize}
    \item \textbf{Age} is the dominant blood pressure predictor (Random Forest importance $\approx 0.67$; OLS $\beta = +0.32$~mmHg systolic per year). Male sex, known blood sugar, and under treatment are associated with higher adjusted blood pressure; self-reported wellbeing is not significant after covariate adjustment.
    \item Individual-level blood pressure prediction remains weak (systolic $R^2 = 0.12$--0.19; diastolic $R^2 = 0.02$--0.10), with mean absolute errors of 10--14~mmHg.
    \item Regional GDP does not significantly predict state-level mean blood pressure. GDP is strongly negatively associated with overweight prevalence at NUTS level~2 ($r = -0.901$). Neighbourhood mean systolic pressure enters neighbourhood effect averaging problem (NEAP) models with a positive coefficient, but the intraclass correlation coefficient (ICC, 1.8\%) indicates that most variance remains at the individual level.
\end{itemize}

\section{Report Structure}
\label{sec:report_structure}

The remainder of this report is organised as follows:
\begin{itemize}
    \item Chapter~\ref{sec:data_analysis} presents exploratory data analysis (Section~\ref{sec:eda}).
    \item Chapter~\ref{sec:modelling} describes the Random Forest modelling pipeline, performance comparisons, and classification extension.
    \item Chapter~\ref{sec:spatial_analysis} reports regional and neighbourhood-level spatial analysis.
    \item Chapter~\ref{sec:conclusion} summarises the principal findings, limitations, and implications.
    \item Appendix~\ref{sec:appendix} contains the artificial intelligence usage declaration.
\end{itemize}

Detailed statistical results and interpretation follow in the subsequent chapters; this introduction establishes context, data characteristics, and the analytical framework only.
